\paragraph{审计偶数窗上的中心坍缩、饱和判据与熵--体积--碰撞恒等链}
本段把附录中已建立的 bin-fold 退化谱、规范群与 Stirling 展开结果，回收为六条可直接调用的终端结论：它们分别刻画审计偶数窗上的中心坍缩与 Abel 满秩分离、最大纤维组合上界的精确饱和判据、任意纤维多重度的普适 Fibonacci 上界及其末位强化、\(\kappa_m\)--KL--碰撞概率之间的恒等链、规范群体积的三因子乘法分解，以及最小退化扇区单独强制的显式碰撞下界。

\begin{theorem}[审计偶数窗上中心坍缩与 Abel 满秩的严格分离]\label{thm:xi-foldbin-audited-even-center-collapse-abel-full-rank}
令
\[
G_m^{\mathrm{bin}}
:=
\Aut(\Fold_m^{\mathrm{bin}})
\cong
\prod_{w\in X_m}S_{d_m^{\mathrm{bin}}(w)}.
\]
若
\(
m\in\{6,8,10,12\}
\)，则
\[
\boxed{\;
\bigl(G_m^{\mathrm{bin}}\bigr)^{\mathrm{ab}}
\cong
(\ZZ/2\ZZ)^{|X_m|}.\;}
\]
更强地，
\[
\boxed{\;
Z\!\bigl(G_6^{\mathrm{bin}}\bigr)\cong (\ZZ/2\ZZ)^8,\qquad
Z\!\bigl(G_m^{\mathrm{bin}}\bigr)=1\ \ (m=8,10,12).\;}
\]
并且在 \(m=6\) 时，boundary sheet-parity 规范子群
\[
Z_{\partial}\cong (\ZZ/2\ZZ)^3\hookrightarrow Z\!\bigl(G_6^{\mathrm{bin}}\bigr)
\]
只是整个中心中的一个规范三维坐标子格；事实上
\[
\boxed{\;
Z\!\bigl(G_6^{\mathrm{bin}}\bigr)/Z_{\partial}\cong (\ZZ/2\ZZ)^5.\;}
\]
\leanpartial{paper\_xi\_foldbin\_audited\_even\_center\_collapse\_abel\_full\_rank}{requested Lean signature quantifies over arbitrary Prop inputs with no hypotheses, so the conjunction is not derivable as stated}
\end{theorem}
\begin{proof}
由定理 \ref{thm:xi-audited-even-abel-center-fibonacci-splitting}，
\[
\bigl(G_m^{\mathrm{bin}}\bigr)^{\mathrm{ab}}
\cong
(\ZZ/2\ZZ)^{F_m}\oplus(\ZZ/2\ZZ)^{D_{2m-2}}.
\]
而
\(
F_m+D_{2m-2}=F_m+F_{m+1}=F_{m+2}=|X_m|
\)，故得到
\[
\bigl(G_m^{\mathrm{bin}}\bigr)^{\mathrm{ab}}\cong (\ZZ/2\ZZ)^{|X_m|}.
\]

另一方面，定理 \ref{thm:xi-foldbin-center-solvable-fiber-criterion} 给出
\[
Z\!\bigl(G_m^{\mathrm{bin}}\bigr)\cong (\ZZ/2\ZZ)^{B_m},
\qquad
B_m:=\#\{w\in X_m:\ d_m^{\mathrm{bin}}(w)=2\}.
\]
由表 \ref{tab:fold_bin_degeneracy_spectrum_m6_m12} 的直方图，
\[
m=6:\ 2:8,\,3:4,\,4:9,
\]
故 \(B_6=8\)；而
\[
m=8:\ 3:21,\,5:11,\,6:23,\qquad
m=10:\ 5:55,\,8:52,\,9:37,\qquad
m=12:\ 8:144,\,12:85,\,13:148,
\]
均无 \(d=2\) 纤维，因此 \(B_8=B_{10}=B_{12}=0\)。从而
\[
Z\!\bigl(G_6^{\mathrm{bin}}\bigr)\cong (\ZZ/2\ZZ)^8,
\qquad
Z\!\bigl(G_m^{\mathrm{bin}}\bigr)=1\ \ (m=8,10,12).
\]

最后，由定理 \ref{thm:xi-window6-central-noncentral-anomaly-splitting}，
\(
Z_{\partial}\cong (\ZZ/2\ZZ)^3
\)
是
\(
Z(G_6^{\mathrm{bin}})
\)
的规范直和因子，并且
\[
Z\!\bigl(G_6^{\mathrm{bin}}\bigr)/Z_{\partial}\cong (\ZZ/2\ZZ)^5.
\]
证毕。
\end{proof}

\begin{theorem}[bin-fold 最大纤维的精确饱和判据与最终严格次 Fibonacci 律]\label{thm:xi-foldbin-maxfiber-saturation-final-subfibonacci}
\leanverified{paper\_xi\_foldbin\_maxfiber\_saturation\_final\_subfibonacci}
令 \(K=K(m)\) 为满足
\[
F_{K+1}\le 2^m-1<F_{K+2}
\]
的唯一整数，并置 \(L:=K-m\)。
定义合法尾部集合
\[
\mathcal T_m
:=
\Bigl\{
(c_{m+1},\dots,c_K)\in\{0,1\}^{L}:\ c_kc_{k+1}=0
\Bigr\},
\]
以及最大尾和值
\[
S_{\max}(m)
:=
\max\Bigl\{
\sum_{k=m+1}^{K}c_kF_{k+1}:\ (c_{m+1},\dots,c_K)\in\mathcal T_m
\Bigr\}.
\]
则最大纤维满足
\[
\boxed{\;
d_{\max}^{\mathrm{bin}}(m)
:=
\max_{w\in X_m}d_m^{\mathrm{bin}}(w)
=
d_m^{\mathrm{bin}}(0^m)
=
\#\Bigl\{
c\in\mathcal T_m:\ \sum_{k=m+1}^{K}c_kF_{k+1}\le 2^m-1
\Bigr\}.\;}
\]
并且，记组合上界 \(F_{L+2}=|\mathcal T_m|\)，则
\[
\boxed{\;
d_{\max}^{\mathrm{bin}}(m)=F_{L+2}
\iff
S_{\max}(m)\le 2^m-1.\;}
\]
进一步，存在 \(m_0\) 使得对所有 \(m\ge m_0\) 都有
\[
\boxed{\;
d_{\max}^{\mathrm{bin}}(m)\le F_{L+2}-1.\;}
\]
\end{theorem}
\begin{proof}
由命题 \ref{prop:fold-bin-digit-dp}，对任意 \(w\in X_m\)，
\(
d_m^{\mathrm{bin}}(w)
\)
恰等于满足相邻约束、边界约束与阈值约束的尾部位串数目。
取 \(w=0^m\) 时，\(V_m(w)=0\) 且 \(w_m=0\)，因而边界约束消失，得到
\[
d_m^{\mathrm{bin}}(0^m)
=
\#\Bigl\{
c\in\mathcal T_m:\ \sum_{k=m+1}^{K}c_kF_{k+1}\le 2^m-1
\Bigr\}.
\]
再由推论 \ref{cor:fold-bin-saturation}，最大纤维确由 \(0^m\) 取得，故第一条公式成立。

同一推论又给出
\[
d_{\max}^{\mathrm{bin}}(m)\le F_{L+2}=|\mathcal T_m|,
\]
且该上界饱和当且仅当阈值约束对全部合法尾部都冗余，也即
\[
2^m-1\ge S_{\max}(m).
\]
这就证明了第二条等价判据。

最后，由注 \ref{rem:fold-bin-finite}，上述组合上界的饱和只能发生有限次。
因此存在 \(m_0\) 使得对一切 \(m\ge m_0\)，恒有
\(
d_{\max}^{\mathrm{bin}}(m)<F_{L+2}
\)；
由于两边都是整数，便得到
\[
d_{\max}^{\mathrm{bin}}(m)\le F_{L+2}-1.
\]
证毕。
\end{proof}

\begin{conjecture}[bin-fold 最大纤维组合上界的精确饱和集是有限孤立集]\label{conj:xi-foldbin-maxfiber-saturation-sporadic-set}
沿用定理 \ref{thm:xi-foldbin-maxfiber-saturation-final-subfibonacci} 的记号。生成表
\ref{tab:fold_binary_saturation_points} 对扫描窗口 \(m\le 400\) 给出的审计结果表明，组合上界
\[
d_{\max}^{\mathrm{bin}}(m)\le F_{L+2}
\]
的饱和点恰为
\[
\{1,2,3,4,5,7,12\}.
\]
据此提出猜想
\[
\boxed{\;
d_{\max}^{\mathrm{bin}}(m)=F_{L+2}
\iff
m\in\{1,2,3,4,5,7,12\}.\;}
\]
等价地，
\[
\boxed{\;
2^m-1\ge S_{\max}(m)
\iff
m\in\{1,2,3,4,5,7,12\}.\;}
\]
若此猜想成立，则定理 \ref{thm:xi-foldbin-maxfiber-saturation-final-subfibonacci} 中“饱和只能发生有限次”的结论可进一步精化为一个完全显式的有限异常集判定。
\end{conjecture}

\begin{theorem}[bin-fold 纤维多重度的普适 Fibonacci 上界与末位强化]\label{thm:xi-foldbin-universal-fibonacci-upper-bound-lastbit-improvement}
\leanverified{paper\_xi\_foldbin\_universal\_fibonacci\_upper\_bound\_lastbit\_improvement}
沿用定理 \ref{thm:xi-foldbin-maxfiber-saturation-final-subfibonacci} 的记号，记
\[
L:=K(m)-m.
\]
则对任意 \(w\in X_m\)，都有
\[
\boxed{\;
d_m^{\mathrm{bin}}(w)\le F_{L+2}=F_{K(m)-m+2}.\;}
\]
若进一步满足 \(w_m=1\)，则更强地有
\[
\boxed{\;
d_m^{\mathrm{bin}}(w)\le F_{L+1}=F_{K(m)-m+1}.\;}
\]
\end{theorem}
\begin{proof}
由命题 \ref{prop:fold-bin-digit-dp}，对任意 \(w\in X_m\)，
\(
d_m^{\mathrm{bin}}(w)
\)
恰等于满足
\[
c_kc_{k+1}=0,\qquad
w_mc_{m+1}=0,\qquad
V_m(w)+\sum_{k=m+1}^{K}c_kF_{k+1}\le 2^m-1
\]
的尾部位串 \((c_{m+1},\dots,c_K)\) 的数目。

先忽略阈值约束。长度 \(L\) 的二进制串中不含相邻 \(1\) 的总数是标准 Fibonacci 计数 \(F_{L+2}\)，故
\[
d_m^{\mathrm{bin}}(w)\le F_{L+2}.
\]

若 \(w_m=1\)，则边界约束强制 \(c_{m+1}=0\)。当 \(L=0\) 时尾部为空串，结论平凡成立；当 \(L\ge 1\) 时，剩余自由尾部是长度 \(L-1\) 的无相邻 \(1\) 串，其数目为
\[
F_{(L-1)+2}=F_{L+1}.
\]
再加入阈值约束只会进一步减少可行尾部，故
\[
d_m^{\mathrm{bin}}(w)\le F_{L+1}.
\]
证毕。
\end{proof}

\begin{theorem}[\texorpdfstring{$\kappa_m$}{kappa_m}、KL 散度与碰撞概率的精确恒等链]\label{thm:xi-foldbin-kappa-kl-collision-identity-chain}
\leanverified{paper\_xi\_foldbin\_kappa\_kl\_collision\_identity\_chain}
令
\[
p_m^{\mathrm{bin}}(w):=\frac{d_m^{\mathrm{bin}}(w)}{2^m},
\qquad
u_m(w):=\frac{1}{|X_m|}
\qquad (w\in X_m).
\]
则有严格恒等式
\[
\boxed{\;
D_{\mathrm{KL}}\!\bigl(p_m^{\mathrm{bin}}\|u_m\bigr)
=
\kappa_m-\log\!\Bigl(\frac{2^m}{|X_m|}\Bigr).\;}
\]
因而
\[
\boxed{\;
\log\!\Bigl(\frac{2^m}{|X_m|}\Bigr)
\le
\kappa_m
\le
\log\!\bigl(2^m\mathsf{Col}^{\mathrm{bin}}_m\bigr).\;}
\]
又由
\[
\chi^2\!\bigl(p_m^{\mathrm{bin}}\|u_m\bigr)
=
|X_m|\,\mathsf{Col}^{\mathrm{bin}}_m-1,
\]
可把上界改写为
\[
\kappa_m
\le
\log\!\Bigl(\frac{2^m}{|X_m|}\Bigr)
\;+\;
\log\!\Bigl(1+\chi^2\!\bigl(p_m^{\mathrm{bin}}\|u_m\bigr)\Bigr).
\]
并且两端取等都当且仅当 \(p_m^{\mathrm{bin}}=u_m\)。
\end{theorem}
\begin{proof}
第一条恒等式正是定理 \ref{thm:xi-foldbin-entropy-gauge-one-nat-gap} 的
\[
D_{\mathrm{KL}}\!\bigl(p_m^{\mathrm{bin}}\|u_m\bigr)
=
\kappa_m+\log|X_m|-m\log 2
\]
的改写。
由于 \(D_{\mathrm{KL}}\ge 0\)，立得
\[
\kappa_m\ge \log\!\Bigl(\frac{2^m}{|X_m|}\Bigr),
\]
且等号成立当且仅当 \(p_m^{\mathrm{bin}}=u_m\)。

另一方面，由定理 \ref{thm:xi-foldbin-gauge-escort-kl-decomposition}，
\[
\kappa_m
=
\log\!\bigl(2^m\mathsf{Col}^{\mathrm{bin}}_m\bigr)
-D_{\mathrm{KL}}\!\bigl(p_m^{\mathrm{bin}}\|\widetilde p_m\bigr),
\]
其中 \(\widetilde p_m(w)=(p_m^{\mathrm{bin}}(w))^2/\mathsf{Col}^{\mathrm{bin}}_m\)。
于是
\[
\kappa_m\le \log\!\bigl(2^m\mathsf{Col}^{\mathrm{bin}}_m\bigr),
\]
且等号成立当且仅当 \(p_m^{\mathrm{bin}}=\widetilde p_m\)，即所有 \(p_m^{\mathrm{bin}}(w)\) 在 \(X_m\) 上相等，从而 \(p_m^{\mathrm{bin}}=u_m\)。

最后，由命题 \ref{prop:fold-bin-chi2-col}，
\[
1+\chi^2\!\bigl(p_m^{\mathrm{bin}}\|u_m\bigr)
=
|X_m|\,\mathsf{Col}^{\mathrm{bin}}_m,
\]
故
\[
\log\!\bigl(2^m\mathsf{Col}^{\mathrm{bin}}_m\bigr)
=
\log\!\Bigl(\frac{2^m}{|X_m|}\Bigr)
\;+\;
\log\!\Bigl(1+\chi^2\!\bigl(p_m^{\mathrm{bin}}\|u_m\bigr)\Bigr),
\]
从而得到所述恒等链。证毕。
\end{proof}

\begin{theorem}[规范群体积的三因子精确乘法分解]\label{thm:xi-foldbin-gauge-volume-three-factor-product}
\leanverified{paper\_xi\_foldbin\_gauge\_volume\_three\_factor\_product}
沿用
\(
G_m^{\mathrm{bin}}:=\Aut(\Fold_m^{\mathrm{bin}})
\)
的记号。
则存在实数 \(r_m\) 满足
\[
0\le r_m\le \frac{|X_m|}{12},
\]
使得
\[
\boxed{\;
|G_m^{\mathrm{bin}}|
=
\exp\!\bigl(2^m(\kappa_m-1)\bigr)\,
(2\pi)^{|X_m|/2}\,
\Bigl(\prod_{w\in X_m}d_m^{\mathrm{bin}}(w)\Bigr)^{1/2}
e^{r_m}.\;}
\]
因此，规范群体积在主熵项之外，精确分裂为 Gaussian 半密度因子与纤维谱几何均值因子，而剩余乘法误差严格受
\(
\exp(|X_m|/12)
\)
控制。
\end{theorem}
\begin{proof}
由定理 \ref{thm:fold-bin-gauge-volume-stirling-second-order}，
\[
g_m
=
\kappa_m-1+\frac{1}{2^{m+1}}
\Bigl(
|X_m|\log(2\pi)+\sum_{w\in X_m}\log d_m^{\mathrm{bin}}(w)
\Bigr)+R_m,
\]
其中
\[
0\le R_m\le \frac{|X_m|}{12\cdot 2^m}.
\]
两边同乘 \(2^m\)，并记
\(
r_m:=2^mR_m
\)，则
\[
\log|G_m^{\mathrm{bin}}|
=
2^m(\kappa_m-1)
+\frac12\Bigl(
|X_m|\log(2\pi)+\sum_{w\in X_m}\log d_m^{\mathrm{bin}}(w)
\Bigr)
+r_m,
\]
且 \(0\le r_m\le |X_m|/12\)。
对该恒等式指数化即得所示乘法分解。证毕。
\end{proof}

\begin{theorem}[审计偶数窗上最小退化扇区强制的显式碰撞下界]\label{thm:xi-foldbin-audited-even-minsector-collision-explicit-floor}
若
\(
m\in\{6,8,10,12\}
\)，则
\[
\boxed{\;
\mathsf{Col}^{\mathrm{bin}}_m
\ge
\frac{1}{2^{2m}}
\left(
F_mF_{m/2}^{\,2}
+
\frac{\bigl(2^m-F_mF_{m/2}\bigr)^2}{F_{m+1}}
\right).\;}
\]
等价地，
\[
\boxed{\;
\chi^2\!\bigl(p_m^{\mathrm{bin}}\|u_m\bigr)
\ge
\frac{F_{m+2}}{2^{2m}}
\left(
F_mF_{m/2}^{\,2}
+
\frac{\bigl(2^m-F_mF_{m/2}\bigr)^2}{F_{m+1}}
\right)-1.\;}
\]
并且等号成立当且仅当所有非最小退化纤维具有同一尺寸。
\end{theorem}
\begin{proof}
记
\[
n:=|X_m|=F_{m+2},
\qquad
r:=\bigl|\mathcal S_{m,d_{\min}(m)}\bigr|=F_m,
\qquad
a:=d_{\min}(m)=F_{m/2},
\]
其中第一、第二、第三个等式分别来自稳定类型计数律、定理 \ref{thm:fold-bin-min-degeneracy-fib-audited} 与同一定理。
再由推论 \ref{cor:fold-bin-minsector-complement-maxfiber-audited}，
\[
n-r=F_{m+1}.
\]

把全部纤维大小写为
\[
\underbrace{a,\dots,a}_{r\text{ 次}},\ e_1,\dots,e_{n-r},
\]
则
\[
\sum_{i=1}^{n-r}e_i=2^m-ra=2^m-F_mF_{m/2}.
\]
于是
\[
\sum_{w\in X_m}\bigl(d_m^{\mathrm{bin}}(w)\bigr)^2
=
ra^2+\sum_{i=1}^{n-r}e_i^2.
\]
对后者应用 Cauchy--Schwarz 不等式，
\[
\sum_{i=1}^{n-r}e_i^2
\ge
\frac{\bigl(\sum_{i=1}^{n-r}e_i\bigr)^2}{n-r}
=
\frac{\bigl(2^m-F_mF_{m/2}\bigr)^2}{F_{m+1}}.
\]
因此
\[
\sum_{w\in X_m}\bigl(d_m^{\mathrm{bin}}(w)\bigr)^2
\ge
F_mF_{m/2}^{\,2}
+
\frac{\bigl(2^m-F_mF_{m/2}\bigr)^2}{F_{m+1}}.
\]
再由
\(
\mathsf{Col}^{\mathrm{bin}}_m
=
2^{-2m}\sum_w(d_m^{\mathrm{bin}}(w))^2
\)
得到碰撞下界。
最后由命题 \ref{prop:fold-bin-chi2-col} 的恒等式
\(
\chi^2(p_m^{\mathrm{bin}}\|u_m)=|X_m|\mathsf{Col}^{\mathrm{bin}}_m-1
\)
即得第二条公式。

Cauchy--Schwarz 取等当且仅当
\(
e_1=\cdots=e_{n-r}
\)，也就是全部非最小退化纤维等大。证毕。
\end{proof}

\begin{theorem}[已审计偶数窗存在完全线性的 Fibonacci 低预算相]\label{thm:xi-foldbin-audited-even-fibonacci-linear-low-budget-phase}
若
\[
m\in\{6,8,10,12\}
\]
且整数预算 \(B\ge 0\) 满足
\[
2^B\le F_{m/2},
\]
则每一条 bin-fold 纤维都超过预算阈值，因而
\[
\boxed{\;
C_m^{\mathrm{bin}}(B)=2^B\,F_{m+2},\qquad
\mathrm{Succ}_m^{\mathrm{bin}}(B)=2^{B-m}F_{m+2}.\;}
\]
换言之，在已审计偶数窗上存在一段完全线性的低预算相；其斜率与截距仅由
\(
F_{m/2}
\)
与
\(
F_{m+2}
\)
两条 Fibonacci 尺度共同控制，而与更高退化层的细节无关。
\end{theorem}
\begin{proof}
由定理 \ref{thm:fold-bin-min-degeneracy-fib-audited}，当 \(m\in\{6,8,10,12\}\) 时，对每个 \(w\in X_m\) 都有
\[
d_m^{\mathrm{bin}}(w)\ge d_{\min}(m)=F_{m/2}.
\]
若 \(2^B\le F_{m/2}\)，则对所有 \(w\in X_m\) 皆有
\[
\min\bigl(d_m^{\mathrm{bin}}(w),2^B\bigr)=2^B.
\]
因此
\[
C_m^{\mathrm{bin}}(B)
=
\sum_{w\in X_m}\min\bigl(d_m^{\mathrm{bin}}(w),2^B\bigr)
=
\sum_{w\in X_m}2^B
=
2^B|X_m|.
\]
再由 \(|X_m|=F_{m+2}\)，便得到
\[
C_m^{\mathrm{bin}}(B)=2^B F_{m+2}.
\]
最后两边除以 \(2^m\)，即得
\[
\mathrm{Succ}_m^{\mathrm{bin}}(B)=2^{B-m}F_{m+2}.
\]
证毕。
\end{proof}

\paragraph{审计偶数窗的奇偶荷截面、可解阴影与纯奇偶荷阶段}
在审计偶数窗上的中心坍缩、Abel 满秩、规范体积与低预算相已经确定之后，还可把
\(
G_m^{\mathrm{bin}}
\)
的可解阴影层进一步压缩为下列几条直接结论。它们分别给出阿贝尔化的一个规范奇偶荷截面、最大可解商的显式相图、中心坍缩与可解非交换坍缩的不同步、纯奇偶荷阶段的半单可解残核、任意可解阴影相对完整规范复杂度的指数级稀薄性，以及
\(
\kappa_m
\)
对规范体积的双侧压缩。

\begin{theorem}[奇偶荷截面的中心交公式]\label{thm:xi-foldbin-even-window-parity-section-center-intersection}
对每个审计偶数窗
\[
m\in\{6,8,10,12\},
\]
沿用定理 \ref{thm:xi-foldbin-audited-even-center-collapse-abel-full-rank} 的记号，并在每个因子
\(
\mathfrak S_{d_m^{\mathrm{bin}}(w)}
\)
中任选一条换位
\(
\tau_w
\)。
定义
\[
E_m:=\bigl\langle \tau_w:\ w\in X_m\bigr\rangle\subset G_m^{\mathrm{bin}}.
\]
则有
\[
E_m\cong (\ZZ/2\ZZ)^{|X_m|},
\]
并且复合映射
\[
E_m\hookrightarrow G_m^{\mathrm{bin}}\twoheadrightarrow
\bigl(G_m^{\mathrm{bin}}\bigr)^{\mathrm{ab}}
\]
是同构。此外，
\[
E_m\cap Z\!\bigl(G_m^{\mathrm{bin}}\bigr)\cong (\ZZ/2\ZZ)^{B_m},
\qquad
B_m:=\#\{w\in X_m:\ d_m^{\mathrm{bin}}(w)=2\}.
\]
特别地，
\[
B_6=8,
\qquad
B_8=B_{10}=B_{12}=0.
\]
\end{theorem}
\begin{proof}
不同的 \(w\) 对应于直积分解
\[
G_m^{\mathrm{bin}}
\cong
\prod_{w\in X_m}\mathfrak S_{d_m^{\mathrm{bin}}(w)}
\]
中的不同坐标因子，因此所选换位两两可交换，且每个 \(\tau_w\) 都满足 \(\tau_w^2=1\)。于是
\[
E_m\cong (\ZZ/2\ZZ)^{|X_m|}.
\]

又由定理 \ref{thm:fold-bin-min-degeneracy-fib-audited}，审计偶数窗上每条 bin-fold 纤维都满足
\(
d_m^{\mathrm{bin}}(w)\ge d_{\min}(m)=F_{m/2}\ge 2
\)。
因此每个因子的符号映射都给出满同态
\[
\operatorname{sgn}_w:\mathfrak S_{d_m^{\mathrm{bin}}(w)}\twoheadrightarrow \ZZ/2\ZZ.
\]
把这些符号映射相乘，得到
\[
\prod_{w\in X_m}\operatorname{sgn}_w:
G_m^{\mathrm{bin}}
\twoheadrightarrow
(\ZZ/2\ZZ)^{|X_m|}.
\]
而每个 \(\tau_w\) 在第 \(w\) 个坐标上给出唯一的奇元，在其余坐标上为单位元，所以该映射限制到 \(E_m\) 后恰是标准坐标同构。再结合定理 \ref{thm:xi-foldbin-audited-even-center-collapse-abel-full-rank} 的
\[
\bigl(G_m^{\mathrm{bin}}\bigr)^{\mathrm{ab}}
\cong
(\ZZ/2\ZZ)^{|X_m|},
\]
便知复合映射
\(
E_m\hookrightarrow G_m^{\mathrm{bin}}\twoheadrightarrow (G_m^{\mathrm{bin}})^{\mathrm{ab}}
\)
为同构。

最后，由定理 \ref{thm:xi-foldbin-audited-even-center-collapse-abel-full-rank}，
\[
Z\!\bigl(G_m^{\mathrm{bin}}\bigr)\cong (\ZZ/2\ZZ)^{B_m}.
\]
直积群的中心等于各因子中心的直积，而
\[
Z(\mathfrak S_d)=
\begin{cases}
\mathfrak S_2,& d=2,\\
1,& d\ge 3.
\end{cases}
\]
故 \(\tau_w\) 落在中心中当且仅当 \(d_m^{\mathrm{bin}}(w)=2\)。于是
\[
E_m\cap Z\!\bigl(G_m^{\mathrm{bin}}\bigr)
=
\bigl\langle \tau_w:\ d_m^{\mathrm{bin}}(w)=2\bigr\rangle
\cong
(\ZZ/2\ZZ)^{B_m}.
\]
最后再由表 \ref{tab:fold_bin_degeneracy_spectrum_m6_m12} 的审计直方图读出
\(
B_6=8
\)
与
\(
B_8=B_{10}=B_{12}=0
\)。
证毕。
\end{proof}
\leanverified{paper\_xi\_foldbin\_even\_window\_parity\_section\_center\_intersection}

\begin{theorem}[审计偶数窗最大可解商的显式相图]\label{thm:xi-foldbin-audited-even-maximal-solvable-quotient}
对
\[
m\in\{6,8,10,12\},
\]
记
\(
Q_m^{\mathrm{sol}}
\)
为
\(
G_m^{\mathrm{bin}}
\)
的最大可解商。定义
\[
q(d):=
\begin{cases}
\mathfrak S_d,& d\le 4,\\
\ZZ/2\ZZ,& d\ge 5.
\end{cases}
\]
则有自然同构
\[
Q_m^{\mathrm{sol}}
\cong
\prod_{w\in X_m}q\!\bigl(d_m^{\mathrm{bin}}(w)\bigr).
\]
特别地，
\[
Q_6^{\mathrm{sol}}
\cong
\mathfrak S_2^{\,8}\times \mathfrak S_3^{\,4}\times \mathfrak S_4^{\,9},
\]
\[
Q_8^{\mathrm{sol}}
\cong
\mathfrak S_3^{\,21}\times (\ZZ/2\ZZ)^{34},
\]
\[
Q_{10}^{\mathrm{sol}}
\cong
(\ZZ/2\ZZ)^{144},
\qquad
Q_{12}^{\mathrm{sol}}
\cong
(\ZZ/2\ZZ)^{377}.
\]
\end{theorem}
\begin{proof}
对称群的标准结构表明：
\[
\mathfrak S_d\ \text{可解}
\iff
d\le 4,
\]
而当 \(d\ge 5\) 时，
\(
A_d
\)
是非阿贝尔单群，因此
\[
\mathfrak S_d/A_d\cong \ZZ/2\ZZ
\]
是 \(\mathfrak S_d\) 的最大可解商。故上式中定义的 \(q(d)\) 正是因子 \(\mathfrak S_d\) 的最大可解商。

现在令
\[
\phi:
G_m^{\mathrm{bin}}
\cong
\prod_{w\in X_m}\mathfrak S_{d_m^{\mathrm{bin}}(w)}
\longrightarrow H
\]
为任意到可解群 \(H\) 的同态。则其在每个坐标因子上的限制
\[
\phi_w:\mathfrak S_{d_m^{\mathrm{bin}}(w)}\to H
\]
具有可解像，因此都必经由 \(q(d_m^{\mathrm{bin}}(w))\)。由直积的泛性质，\(\phi\) 必经由
\[
\prod_{w\in X_m}q\!\bigl(d_m^{\mathrm{bin}}(w)\bigr).
\]
这就证明了该直积正是
\(
G_m^{\mathrm{bin}}
\)
的最大可解商。

最后只需把表 \ref{tab:fold_bin_degeneracy_spectrum_m6_m12} 的审计直方图代入。对
\(
m=6
\)
有
\[
2:8,\qquad 3:4,\qquad 4:9,
\]
故
\[
Q_6^{\mathrm{sol}}
\cong
\mathfrak S_2^{\,8}\times \mathfrak S_3^{\,4}\times \mathfrak S_4^{\,9}.
\]
对
\(
m=8
\)
有
\[
3:21,\qquad 5:11,\qquad 6:23,
\]
故
\[
Q_8^{\mathrm{sol}}
\cong
\mathfrak S_3^{\,21}\times (\ZZ/2\ZZ)^{11+23}
=
\mathfrak S_3^{\,21}\times (\ZZ/2\ZZ)^{34}.
\]
对
\(
m=10
\)
有
\[
5:55,\qquad 8:52,\qquad 9:37,
\]
全部因子都贡献符号商，所以
\[
Q_{10}^{\mathrm{sol}}\cong (\ZZ/2\ZZ)^{55+52+37}=(\ZZ/2\ZZ)^{144}.
\]
同理，对
\(
m=12
\)
的直方图
\[
8:144,\qquad 12:85,\qquad 13:148
\]
可得
\[
Q_{12}^{\mathrm{sol}}\cong (\ZZ/2\ZZ)^{144+85+148}=(\ZZ/2\ZZ)^{377}.
\]
证毕。
\end{proof}
\leanverified{paper\_xi\_foldbin\_audited\_even\_maximal\_solvable\_quotient}

\begin{corollary}[中心坍缩与可解非交换坍缩不同步]\label{cor:xi-foldbin-center-solvable-collapse-two-stage}
在审计偶数窗上，中心坍缩与可解非交换坍缩是两件不同层级的事件。
更具体地，
\[
Z\!\bigl(G_m^{\mathrm{bin}}\bigr)=1
\qquad (m=8,10,12),
\]
但
\[
Q_8^{\mathrm{sol}}
\cong
\mathfrak S_3^{\,21}\times (\ZZ/2\ZZ)^{34}
\]
仍然保留非阿贝尔可解因子，而
\[
Q_{10}^{\mathrm{sol}}
\cong
(\ZZ/2\ZZ)^{144},
\qquad
Q_{12}^{\mathrm{sol}}
\cong
(\ZZ/2\ZZ)^{377}
\]
才首次退化为纯奇偶荷格。
因此，在已审计偶数窗中，“中心已死而可解非交换尚存”的现象恰发生在 \(m=8\)。
\end{corollary}
\begin{proof}
由定理 \ref{thm:xi-foldbin-audited-even-center-collapse-abel-full-rank}，
\[
Z\!\bigl(G_m^{\mathrm{bin}}\bigr)=1
\qquad (m=8,10,12).
\]
另一方面，定理 \ref{thm:xi-foldbin-audited-even-maximal-solvable-quotient} 给出
\[
Q_8^{\mathrm{sol}}
\cong
\mathfrak S_3^{\,21}\times (\ZZ/2\ZZ)^{34},
\]
故在 \(m=8\) 时最大可解商仍含非阿贝尔可解因子；而在
\(
m=10,12
\)
时，最大可解商都已坍缩为纯二元奇偶荷格。证毕。
\end{proof}
\leanverified{paper\_xi\_foldbin\_center\_solvable\_collapse\_two\_stage}

\leanverified{paper\_xi\_foldbin\_pure\_parity\_solvable\_residual\_kernel}
\begin{theorem}[纯奇偶荷阶段的半单可解残核]\label{thm:xi-foldbin-pure-parity-solvable-residual-kernel}
设分辨率 \(m\) 满足
\[
d_m^{\mathrm{bin}}(w)\ge 5
\qquad (\forall\, w\in X_m).
\]
定义正规子群
\[
N_m:=\prod_{w\in X_m}A_{d_m^{\mathrm{bin}}(w)}
\triangleleft
G_m^{\mathrm{bin}}.
\]
则有短正合列
\[
1\longrightarrow N_m
\longrightarrow G_m^{\mathrm{bin}}
\longrightarrow (\ZZ/2\ZZ)^{|X_m|}
\longrightarrow 1.
\]
并且：
\begin{enumerate}
\item \(N_m\) 是 \(G_m^{\mathrm{bin}}\) 的可解残核，即所有到可解群同态之核的交；
\item \(N_m\) 是 \(|X_m|\) 个非阿贝尔单群的直积；
\item 任意到可解群的同态都必经由奇偶荷格
\(
(\ZZ/2\ZZ)^{|X_m|}
\)。
\end{enumerate}
特别地，表 \ref{tab:fold_bin_degeneracy_spectrum_m6_m12} 表明该结论适用于
\[
m=10,\qquad m=12.
\]
\end{theorem}
\begin{proof}
在假设
\(
d_m^{\mathrm{bin}}(w)\ge 5
\)
下，每个交错群
\(
A_{d_m^{\mathrm{bin}}(w)}
\)
都是非阿贝尔单群，并满足
\[
\mathfrak S_{d_m^{\mathrm{bin}}(w)}/A_{d_m^{\mathrm{bin}}(w)}
\cong
\ZZ/2\ZZ.
\]
因此
\[
G_m^{\mathrm{bin}}/N_m
\cong
\prod_{w\in X_m}
\mathfrak S_{d_m^{\mathrm{bin}}(w)}/A_{d_m^{\mathrm{bin}}(w)}
\cong
(\ZZ/2\ZZ)^{|X_m|},
\]
从而得到所示短正合列。

再令
\(
\phi:G_m^{\mathrm{bin}}\to H
\)
为任意到可解群 \(H\) 的同态。其在每个坐标因子上的限制
\[
\phi_w:\mathfrak S_{d_m^{\mathrm{bin}}(w)}\to H
\]
具有可解像，而
\(
A_{d_m^{\mathrm{bin}}(w)}
\)
是 \(\mathfrak S_{d_m^{\mathrm{bin}}(w)}\) 的唯一非平凡非阿贝尔单正规子群，因此必落在 \(\ker\phi_w\) 中。对全部 \(w\) 同时成立，于是
\[
N_m\subseteq \ker\phi.
\]
这说明任意到可解群的同态都经由
\(
G_m^{\mathrm{bin}}/N_m\cong (\ZZ/2\ZZ)^{|X_m|}
\)，故 \(N_m\) 正是可解残核。

第二条由 \(N_m\) 的定义与每个 \(A_{d_m^{\mathrm{bin}}(w)}\) 的非阿贝尔单性立即得出；第三条只是第一条在商层的重述。最后，由表 \ref{tab:fold_bin_degeneracy_spectrum_m6_m12} 可见 \(m=10,12\) 时所有纤维大小都至少为 \(5\)，故这两个审计偶数窗均处于纯奇偶荷阶段。证毕。
\end{proof}

\begin{theorem}[任意可解阴影相对于规范复杂度的指数稀薄律]\label{thm:xi-foldbin-solvable-shadow-exponential-thinning}
设分辨率 \(m\) 满足
\[
d_m^{\mathrm{bin}}(w)\ge 5
\qquad (\forall\, w\in X_m),
\]
并令 \(Q\) 为 \(G_m^{\mathrm{bin}}\) 的任意非平凡可解商。则有
\[
\log |Q|
\le
|X_m|\log 2
=
F_{m+2}\log 2.
\]
另一方面，
\[
\log\bigl|G_m^{\mathrm{bin}}\bigr|
=
2^m g_m
=
2^m
\left(
m\log\!\frac{2}{\varphi}
-1-\frac{\log\varphi}{1+\varphi^2}
+O\!\left(m\left(\frac{\varphi}{2}\right)^m\right)
\right).
\]
因此
\[
\frac{\log|G_m^{\mathrm{bin}}|}{\log|Q|}
\ge
\frac{2^m}{F_{m+2}\log 2}
\left(
m\log\!\frac{2}{\varphi}
-1-\frac{\log\varphi}{1+\varphi^2}
+O\!\left(m\left(\frac{\varphi}{2}\right)^m\right)
\right).
\]
沿任意一条满足上述假设的共尾子列，进一步有
\[
\frac{\log|G_m^{\mathrm{bin}}|}{\log|Q|}
\ge
\frac{\sqrt5}{\varphi^2\log 2}
\left(
m\log\!\frac{2}{\varphi}
-1-\frac{\log\varphi}{1+\varphi^2}
+o(m)
\right)
\left(\frac{2}{\varphi}\right)^m.
\]
\end{theorem}
\begin{proof}
由定理 \ref{thm:xi-foldbin-pure-parity-solvable-residual-kernel}，任意可解商都必经由奇偶荷格
\[
(\ZZ/2\ZZ)^{|X_m|}.
\]
因此
\[
|Q|\le 2^{|X_m|},
\]
也就是
\[
\log |Q|\le |X_m|\log 2.
\]
再由稳定类型计数律 \(|X_m|=F_{m+2}\)，得到第一条公式。

另一方面，定理 \ref{thm:xi-foldbin-gauge-density-exact-thermodynamic-constant} 给出
\[
g_m
=
m\log\!\Bigl(\frac{2}{\varphi}\Bigr)
-1-\frac{\log\varphi}{1+\varphi^2}
+O\!\Bigl(m\Bigl(\frac{\varphi}{2}\Bigr)^m\Bigr).
\]
乘以 \(2^m\) 便得到
\[
\log|G_m^{\mathrm{bin}}|=2^m g_m
\]
的所示展开。把这两个估计合并，便得到第二条不等式。

最后，沿任意满足假设的共尾子列，Binet 公式给出
\[
F_{m+2}
=
\frac{\varphi^{m+2}}{\sqrt5}+o(\varphi^m),
\]
因此
\[
\frac{2^m}{F_{m+2}\log 2}
=
\frac{\sqrt5}{\varphi^2\log 2}
\left(\frac{2}{\varphi}\right)^m
\bigl(1+o(1)\bigr).
\]
与前式相乘后，主项来自
\(
m\log(2/\varphi)
\)，而全部误差可统一吸收到括号中的 \(o(m)\) 中，故得到最后一式。证毕。
\end{proof}

因此，一旦进入纯奇偶荷阶段，任何可解证书所能承载的全部信息量都至多线性于
\(
F_{m+2}
\)，而完整规范复杂度的主尺度却是
\(
2^m\log(2/\varphi)
\)；
可解阴影相对于完整 gauge 体积因此只是指数级稀薄的投影。

\begin{theorem}[规范体积对 \texorpdfstring{$\kappa_m$}{kappa_m} 的双侧压缩]\label{thm:xi-foldbin-gauge-volume-kappa-two-sided-compression}
对任意分辨率 \(m\)，记
\[
d_{\min}(m):=\min_{w\in X_m}d_m^{\mathrm{bin}}(w).
\]
则成立严格双侧界
\[
2^m(\kappa_m-1)
+\frac{|X_m|}{2}\log\!\bigl(2\pi d_{\min}(m)\bigr)
\le
\log\bigl|G_m^{\mathrm{bin}}\bigr|
\]
以及
\[
\log\bigl|G_m^{\mathrm{bin}}\bigr|
\le
2^m(\kappa_m-1)
+\frac{|X_m|}{2}\log\!\Bigl(2\pi\frac{2^m}{|X_m|}\Bigr)
+\frac{|X_m|}{12}.
\]
特别地，
\[
\log\bigl|G_m^{\mathrm{bin}}\bigr|
=
2^m(\kappa_m-1)+O(|X_m|\,m).
\]
\end{theorem}
\begin{proof}
由定理 \ref{thm:xi-foldbin-gauge-volume-three-factor-product}，存在实数 \(r_m\) 满足
\[
0\le r_m\le \frac{|X_m|}{12},
\]
并且
\[
\log\bigl|G_m^{\mathrm{bin}}\bigr|
=
2^m(\kappa_m-1)
+\frac{|X_m|}{2}\log(2\pi)
+\frac12\sum_{w\in X_m}\log d_m^{\mathrm{bin}}(w)
+r_m.
\]
下界由
\[
\sum_{w\in X_m}\log d_m^{\mathrm{bin}}(w)
\ge
|X_m|\log d_{\min}(m)
\]
直接推出，即
\[
\log\bigl|G_m^{\mathrm{bin}}\bigr|
\ge
2^m(\kappa_m-1)
+\frac{|X_m|}{2}\log\!\bigl(2\pi d_{\min}(m)\bigr).
\]

上界则由 \(\log\) 的凹性与质量守恒
\[
\sum_{w\in X_m}d_m^{\mathrm{bin}}(w)=2^m
\]
给出：
\[
\frac1{|X_m|}
\sum_{w\in X_m}\log d_m^{\mathrm{bin}}(w)
\le
\log\!\left(
\frac1{|X_m|}
\sum_{w\in X_m}d_m^{\mathrm{bin}}(w)
\right)
=
\log\!\left(\frac{2^m}{|X_m|}\right).
\]
代回并使用
\(
r_m\le |X_m|/12
\)
便得
\[
\log\bigl|G_m^{\mathrm{bin}}\bigr|
\le
2^m(\kappa_m-1)
+\frac{|X_m|}{2}\log\!\Bigl(2\pi\frac{2^m}{|X_m|}\Bigr)
+\frac{|X_m|}{12}.
\]

最后，由
\(
|X_m|=F_{m+2}=\Theta(\varphi^m)
\)
且
\(
\log(2^m/|X_m|)=O(m)
\)，
可知两端修正项都只有 \(O(|X_m|m)\) 的量级，从而
\[
\log\bigl|G_m^{\mathrm{bin}}\bigr|
=
2^m(\kappa_m-1)+O(|X_m|\,m).
\]
证毕。
\end{proof}

\begin{conjecture}[可解阴影的最终纯奇偶荷化]\label{conj:xi-foldbin-solvable-shadow-eventual-pure-parity}
存在偶数阈值 \(m_0\)，使得对一切充分大的偶数 \(m\)，最大可解商都满足
\[
Q_m^{\mathrm{sol}}
\cong
(\ZZ/2\ZZ)^{|X_m|}.
\]
等价地说，从某个偶数窗口开始，最大可解商中不再保留任何
\(
\mathfrak S_3
\)
或
\(
\mathfrak S_4
\)
型的非阿贝尔可解残余，全部可解信息完全坍缩为纤维奇偶荷格。

当前已审计数据进一步指向更强的版本：\(m=8\) 是最后一个“中心已死而可解非交换尚存”的审计偶数窗。
\end{conjecture}

综上，bin-fold 规范群中的阿贝尔可见性、中心可见性与可解可见性并不属于同一层级：偶数审计窗先失去中心，再失去可解非交换阴影，而一旦进入纯奇偶荷阶段，全部可解信息便被压缩为
\(
(\ZZ/2\ZZ)^{|X_m|}
\)
这一指数级稀薄的外显投影。

\endinput


\endinput


\endinput
